\documentclass[UTF8]{ctexart}
\usepackage{xeCJK}
\usepackage{geometry}
\geometry{left=2cm,right=2cm,top=2.5cm,bottom=2.5cm}
\setlength{\headheight}{12.7pt}

\usepackage{titlesec}
\titleformat{\section}
  {\normalfont\large\bfseries}
  {\thesection}
  {1em}
  {}
\titlespacing*{\section}
  {0pt}
  {3.5ex plus 1ex minus .2ex}
  {2.3ex plus .2ex}

\usepackage{amsmath}
\usepackage{amssymb}
\usepackage{amsthm}
\usepackage{enumitem}
\setlist[enumerate]{itemsep=0pt,topsep=0pt,parsep=0pt,leftmargin=0pt,itemindent=2.5em}
\setlist[itemize]{itemsep=0pt,topsep=0pt,parsep=0pt,leftmargin=0pt,itemindent=2.5em}
\usepackage{fancyhdr}
\pagestyle{fancy}
\fancyhf{}
\usepackage{graphicx}
\usepackage{hyperref}
\usepackage{indentfirst}
\usepackage{lmodern}

\begin{document}
\setlength{\abovedisplayskip}{1pt}
\setlength{\belowdisplayskip}{1pt}

\section{基数指数}

\hypertarget{sec:定义1.1}{}
\noindent\textbf{定义1.1 (基数指数)}

对任意的\href{../01 基数/01 基数#nameddest=sec:定义1.1}{基数}
$\kappa,\lambda$, 定义
\[
    \kappa^\lambda\overset{\mathrm{def}}{=}|\,^\lambda\kappa\,|,
\]
称之为\textbf{基数指数}.

\noindent{\kaishu 备注.
基数指数与序数指数不同, 这里的符号都表示基数指数, $+$和$\cdot$也表示基数运算.}

\vspace{10pt}

\hypertarget{sec:定理1.1}{}
\noindent\textbf{定理1.1}

任给基数$\kappa,\lambda,\mu$,
\begin{enumerate}
    \item $\kappa^{\lambda+\mu}=\kappa^\lambda\kappa^\mu$,
    \item $\kappa^{\lambda\mu}=(\kappa^\lambda)^\mu$,
    \item 如果$\lambda<\mu$, 那么$\kappa^\lambda\leqslant\kappa^\mu$,
    \item 如果$\lambda<\mu$, 那么$\lambda^\kappa\leqslant\mu^\kappa$.
\end{enumerate}

\noindent{\kaishu 证明.}

\begin{enumerate}
    \item $\kappa^{\lambda+\mu}\approx\,\!^{\lambda+\mu}\kappa\approx
    \,\!^{\lambda\sqcup\mu}\kappa\approx
    \,\!^\lambda\kappa\times\,\!^\mu\kappa\approx\kappa^\lambda\kappa^\mu$,
    \item $\kappa^{\lambda\mu}\approx\,\!^{\lambda\mu}\kappa\approx
    \,\!^{\mu\times\lambda}\kappa\approx
    \,\!^\mu(\,\!^\lambda\kappa)\approx(\kappa^\lambda)^\mu$,
    \item 如果$\lambda<\mu$, 那么$\,\!^\lambda\kappa\preccurlyeq\,\!^\mu\kappa$,
    \item 如果$\lambda<\mu$, 那么$\,\!^\kappa\lambda\subset\,\!^\kappa\mu$.
    \hfill $\qedsymbol$
\end{enumerate}

\vspace{10pt}

当$\lambda$有限时, 归纳即得下面的定理.

\vspace{10pt}

\hypertarget{sec:定理1.2}{}
\noindent\textbf{定理1.2}

任给基数$\kappa,\lambda$,
\begin{enumerate}
    \item 如果$\kappa,\lambda$都有限, 那么$\kappa^\lambda$
    等同于\href{../../03 自然数/02 自然数算术/03 自然数的指数#nameddest=sec:定义1.1}
    {自然数的指数},
    \item 如果$\kappa$无限, $\lambda$有限非零, 那么$\kappa^\lambda=\kappa$.
\end{enumerate}

\vspace{10pt}

\hypertarget{sec:定理1.3}{}
\noindent\textbf{定理1.3}

任给基数$\kappa,\lambda$,
\begin{enumerate}
    \item $2^\lambda=|\mathcal{P}(\lambda)|$,
    \item 如果$\lambda$无限且$2\leqslant\kappa\leqslant\lambda$,
    那么$\kappa^\lambda=2^\lambda$.
\end{enumerate}


\noindent{\kaishu 证明.}

\begin{enumerate}
    \item 首先对任意的$p\in\,\!^\lambda2$定义
    \[
        f(p)=\{x\in\lambda\mid p(x)=0\},
    \]
    显然$f:\,\!^\lambda2\to\mathcal{P}(\lambda)$是双射,
    从而$2^\lambda=|\mathcal{P}(\lambda)|$.
    \item 如果$\lambda$无限且$2\leqslant\kappa\leqslant\lambda$, 那么
    \[
        \kappa^\lambda\leqslant\lambda^\lambda\approx\,\!^\lambda\lambda
        \preccurlyeq\mathcal{P}(\lambda\times\lambda)\approx
        \mathcal{P}(\lambda)\approx 2^\lambda,
    \]
    同时$\kappa^\lambda\geqslant 2^\lambda$,
    所以$\kappa^\lambda=2^\lambda$. \hfill $\qedsymbol$
\end{enumerate}

% 基数就先到这里, 大概够用了.

\end{document}
